\documentclass[10pt,twocolumn,a4paper]{article}

% -------------------------------------------------
% Idioma y codificación
% -------------------------------------------------
\usepackage[T1]{fontenc}
\usepackage{lmodern}
\usepackage[english]{babel}

% -------------------------------------------------
% Formato del documento
% -------------------------------------------------
\usepackage[
    top=1.8cm,
    bottom=1.8cm,
    left=1.8cm,
    right=1.8cm,
    columnsep=0.7cm
]{geometry}

\usepackage{microtype}
\usepackage{parskip}

% -------------------------------------------------
% Matemáticas
% -------------------------------------------------
\usepackage{amsmath}
\usepackage{amssymb}
\usepackage{bm}

% -------------------------------------------------
% Figuras y tablas
% -------------------------------------------------
\usepackage{graphicx}
\usepackage{booktabs}
\usepackage{siunitx}

% -------------------------------------------------
% Enlaces y referencias
% -------------------------------------------------
\usepackage[hidelinks]{hyperref}

% Carpeta de imágenes
\graphicspath{{images/}}

% -------------------------------------------------
% Información del documento
% -------------------------------------------------
\title{
    \textbf{Angular Goal Analysis\\
    and\\
    Robot Heading Error}
}

\author{
    Julio Quejido Barrera\\
    Humanoide Robotik\\
    Julio's Laboratories
}

\date{\today}

% =================================================
\begin{document}
% =================================================

\maketitle

\begin{abstract}
This project extends the geometric goal analysis developed 
in the previous projects by introducing an angular representation 
of the relative orientation between a mobile robot and multiple goals. 
For each candidate goal, the absolute direction of the corresponding 
robot-to-goal vector is calculated using the two-argument arctangent 
function $\operatorname{atan2}$. The difference between this goal angle 
and the robot orientation produces an angular error. Because angular 
quantities are periodic, the raw error is normalized to obtain the 
shortest signed rotational difference within the interval $[-\pi,\pi)$. 
The sign of the normalized angular error is then used to determine 
whether the robot should rotate left, rotate right, or continue moving 
straight.
\end{abstract}

\section{Introduction}

A mobile robot may operate in an environment containing 
several possible target positions. In this situation, 
selecting a goal cannot depend only on the distance 
between the robot and each candidate. The robot must 
also consider its current orientation, because a nearby 
goal located behind the robot may be less relevant than 
a goal positioned within its forward field.

This project introduces a robot heading vector and uses 
it to evaluate the relative position of multiple candidate 
goals. Each goal is processed individually, classified 
according to its position with respect to the robot heading, 
and tested using a visibility condition.

The geometric operations developed in Project 1 are reused 
to obtain the distance and normalized direction associated 
with each candidate goal. The new contribution of this project 
is therefore the combination of robot orientation, scalar-product 
analysis, spatial classification and visibility filtering.

\section{Theoretical Background}

\subsection{Angular Representation of a Vector}

For goal $i$, let $\mathbf{v}_{RG_i}$ denote the robot-to-goal vector. A two-dimensional direction vector can also be represented by an angle measured from the positive global $x$-axis. The absolute angle toward goal $i$ is

\begin{equation}
    \theta_{G_i} = \operatorname{atan2}(v_{y_i}, v_{x_i}).
    \label{eq:goal_angle}
\end{equation}

The value $\theta_{G_i}$ describes the orientation of the robot-to-goal vector in the global coordinate system.

The two-argument arctangent is preferred over

\begin{equation}
    \arctan\left(\frac{v_{y_i}}{v_{x_i}}\right)
\end{equation}

because the one-argument expression cannot uniquely determine the 
quadrant of the vector and becomes undefined when $v_{x_i}=0$. 
In contrast, $\operatorname{atan2}$ considers the signs of both 
vector components and returns the correct angular quadrant.


\begin{table}[htbp]
\centering
\caption{Comparison Between $\arctan(y/x)$ and $\text{atan2}(y, x)$.}
\label{tab:atan2_comparison}
\small % Reduces font size slightly to guarantee it fits the column width
\begin{tabular}{lcccc}
\toprule
\textbf{Vec.} $(x, y)$ & \textbf{Quad.} & $\arctan(y/x)$ & $\text{atan2}(y, x)$ & \textbf{Correct} \\
\midrule
$(3, 3)$   & I   & $45^{\circ}$ & $45^{\circ}$   & Yes \\
$(-3, -3)$ & III & $45^{\circ}$ & $-135^{\circ}$ & No \\
\bottomrule
\end{tabular}
\end{table}


\subsection{Absolute Goal Angle and Robot Orientation}

The robot orientation is represented by $\theta_R$, while 
the absolute direction toward goal $i$ is represented by $\theta_{G_i}$. 
Both angles are measured from the same positive global $x$-axis. 
The robot orientation describes the direction in which the robot is currently 
facing, whereas the goal angle describes the global direction from 
the robot position toward the selected target. These quantities 
are absolute angles. The goal angle alone does not indicate how much 
the robot must rotate, because the current robot orientation must also 
be considered. A relative angular difference must therefore be calculated.

\subsection{Angular Error}

The raw angular error for goal $i$ is

\begin{equation}
    e_{\theta_i} = \theta_{G_i} - \theta_R.
    \label{eq:raw_angular_error}
\end{equation}

This value represents the signed rotational difference between the current robot orientation and the direction toward the goal.

Before normalization,

\[
    e_{\theta_i} > 0
\]

indicates a counterclockwise angular difference, while

\[
    e_{\theta_i} < 0
\]

indicates a clockwise angular difference. When

\[
    e_{\theta_i} = 0,
\]

the robot orientation is aligned with the goal direction.

The magnitude

\[
    \left|e_{\theta_i}\right|
\]

indicates the size of the angular difference. However, direct angle 
subtraction does not always provide the shortest physical rotation 
because angular quantities are periodic.

\subsection{Periodicity of Angles}

Angular quantities repeat after every complete revolution. Therefore,

\[
    \theta = \theta + 2k\pi,
    \qquad
    k \in \mathbb{Z}.
\]

Angles that differ by an integer multiple of $2\pi$ represent the same 
physical orientation.

For example, consider

\[
    \theta_R = -179^\circ
\]

and

\[
    \theta_G = 179^\circ.
\]

Direct subtraction produces

\[
    e_\theta
    =
    179^\circ - \left(-179^\circ\right)
    =
    358^\circ.
\]

However, the shortest rotation from the robot orientation to the goal 
direction is

\[
    -2^\circ.
\]

The large raw error appears because the two angles lie on opposite sides 
of the $-\pi/\pi$ boundary. This discontinuity shows that direct 
subtraction alone is not sufficient for obtaining the shortest rotational 
difference.

\subsection{Angular Error Normalization}

Because angular quantities are periodic, the raw angular error may not 
represent the shortest physical rotation between the robot orientation 
and the goal direction. The project therefore normalizes the error using

\begin{equation}
    e_{\theta_i}^{\mathrm{norm}}
    =
    \left(e_{\theta_i} + \pi\right)
    \bmod \left(2\pi\right)
    - \pi.
    \label{eq:normalized_angular_error}
\end{equation}

This operation maps any angular error to the interval

\[
    -\pi \leq e_{\theta_i}^{\mathrm{norm}} < \pi.
\]

The normalized value represents the shortest signed angular difference 
between the current robot orientation and the direction toward goal $i$.

When

\[
    e_{\theta_i}^{\mathrm{norm}} \approx 0,
\]

the robot is approximately aligned with the goal direction.

When

\[
    \left|e_{\theta_i}^{\mathrm{norm}}\right| \approx \pi,
\]

the goal direction is approximately opposite to the robot orientation. 
Because the normalization interval is $[-\pi,\pi)$, an angular difference 
of exactly $\pi$ is represented as $-\pi$.

Small positive or negative values indicate that only a small rotation is 
required. The sign determines the direction of the shortest rotation, 
while the magnitude indicates its size. This normalization is essential 
in robotics because orientation variables wrap around after every 
complete revolution.

\subsection{Rotation Direction Decision}

The normalized angular error is converted into a simple motion decision:

\[
\begin{aligned}
    e_{\theta_i}^{\mathrm{norm}} < 0
    &\quad \Rightarrow \quad
    \text{turn right}, \\
    e_{\theta_i}^{\mathrm{norm}} = 0
    &\quad \Rightarrow \quad
    \text{go straight}, \\
    e_{\theta_i}^{\mathrm{norm}} > 0
    &\quad \Rightarrow \quad
    \text{turn left}.
\end{aligned}
\]

The sign determines the shortest rotation direction, while the magnitude 
indicates how large the required rotation is.

\section{Project Task}

Define one robot position, a robot orientation angle, and several 
candidate goal positions as NumPy arrays. For every goal, compute the 
absolute direction using $\operatorname{atan2}$, calculate the angular 
error between the goal direction and the robot orientation, and normalize 
the error to the interval $[-\pi,\pi)$. Convert the angular values to 
degrees for interpretation and use the sign of the normalized error to 
determine whether the robot should turn left, turn right, or continue 
straight.


% =================================================
\end{document}
% =================================================